\section{Conclusion and Outlook}

The central question behind this project work was to find out whether a differentiable physics simulator, 
used as the prior mean in Bayesian Optimization, can actually help with a realistic accelerator tuning problem(ARES Experimental Area) 
or does it only work on a toy example(FODO lattice)? Based on everything presented, including implementation, experiments and the results, it can confidently be said "yes". \

The problem tackled in this project was the tranverse beam tuning task
at the ARES Experimental Area, the idea of Cheetah based prior demonstrated on a simple two dimensional FODO lattice \cite{kaiser2024bridging} 
was scaled up to a five dimensional beam tuning task now under this project scope, including 5 tuneable magnets (three quadrupole
magnets and two corrector magnets), six unknown quadrupole
misalignment values that the prior had to learn during the optimization and a four components based objective that asks
optimizer to focus and centre the beam simultaneaously. 
The code was built around a module \texttt{AresPriorMean} that
wraps the Cheetah differentiable simulator and plugs into the Xopt framework. The project contains a modular codebase architecture, with
clean separation between the prior mean function, the optimization loop, evaluation and plotting. The approach was tested against three experimental conditions. In the matched scenario,
the prior already knows the correct system state, which is the best-case scenario. 
In the mismatched case, the modified beam and wrong misalignment values were given to prior on purpose, 
forcing it to learn  the correct values during optimization. And in the matched\_prior\_newtask
scenario, the correct misalignment values were given upfront, which was the matched case for evaluation and comparison purposes under this project.
Both Cheetah prior based BO were compared against two baseline techniques, the standard BO
with an uninformed prior and Nelder-Mead. 

The results speak for themselves. Both the physics informed BO variants (matched and mismatched) achieved a median MAE of 6.22$\mu$m and 6.49$\mu$m respectively, which is significantly higher than standard BO (45.88$\mu$m) and Nelder-Mead(67.13$\mu$m).
The matched prior hit the 40$\mu$m target in a median 4 steps, the mismatched prior in 14, whereas standard BO required 82.5 steps on a rare occasions (20\% success rate) and 
Nelder-Mead never reached the target. What makes the approach especially promising for real world use is its robustness. Even when starting with incorrect misalignments values 
and modified beam distribution (the mismatched prior in BO case), the prior still outperformed the uninformed zero mean prior (standard BO) and caught up to the perfectly matched prior scenario 
of BO very soon within 15-20 iterations. The trainable misalignment parameters successfully converged very close to the ground truth purely from optimization data, which means that the method does 
not just tune the beam but it also learns something about the machine’s physical state along the way.  \

Ofcourse, the ARES Experimental Ares with its five tuneable parameters is still a relatively small problem compared to large-scale facilities like European XFEL and PETRA III, 
but this is a significant step forward from the two dimensional FODO lattice that was previously tested to evaluate cheetah’s potential. The fact that the method scales gracefully 
to this level sends quite good signs for scaling further and handling more complex beamlines in future. 



